% !TeX program = xelatex
% Все численные координаты рисунков и таблиц включены в этот файл.
\documentclass[11pt,a4paper]{article}
\usepackage{fontspec}
\usepackage[russian]{babel}
\IfFontExistsTF{Cambria}{\setmainfont{Cambria}}{\setmainfont{DejaVu Serif}}
\IfFontExistsTF{Calibri}{\setsansfont{Calibri}}{\setsansfont{DejaVu Sans}}
\usepackage[a4paper,left=21mm,right=19mm,top=22mm,bottom=20mm,headheight=15pt,headsep=5mm,footskip=10mm]{geometry}
\usepackage{amsmath,booktabs,array,graphicx,xcolor,fancyhdr,hyperref,pgfplots,caption,titlesec}
\pgfplotsset{compat=1.16}
\definecolor{ink}{HTML}{193644}\definecolor{teal}{HTML}{066E70}\definecolor{rust}{HTML}{9E4F39}
\definecolor{plum}{HTML}{73569A}\definecolor{ochre}{HTML}{916316}\definecolor{sky}{HTML}{34659E}
\definecolor{muted}{HTML}{4D626A}\definecolor{hair}{HTML}{CFD9DA}
\hypersetup{unicode=true,colorlinks=true,linkcolor=teal,urlcolor=teal,pdftitle={Ti–Al. Замыкание сольвуса на MTP},pdfsubject={Атомистические расчёты при нулевом давлении}}
\captionsetup{font=small,labelfont={bf,color=teal},labelsep=period,singlelinecheck=false}
\titleformat{\section}{\sffamily\bfseries\Large\color{ink}}{}{0pt}{}
\titleformat{\subsection}{\sffamily\bfseries\large\color{ink}}{}{0pt}{}
\titlespacing*{\section}{0pt}{4pt}{7pt}\titlespacing*{\subsection}{0pt}{9pt}{3pt}
\setlength{\parindent}{0pt}\setlength{\parskip}{4pt}\setlength{\emergencystretch}{2em}
\renewcommand{\arraystretch}{1.18}\raggedbottom\pagestyle{fancy}\fancyhf{}
\fancyhead[L]{\small\sffamily\color{muted}Ti--Al\quad/\quad Равновесие на ГПУ-решётке}
\fancyhead[R]{\small\sffamily\color{muted}Прямой MTP\quad/\quad \(p=0\)}
\fancyfoot[R]{\small\sffamily\color{muted}\thepage}
\renewcommand{\headrulewidth}{.25pt}
\begin{document}
\thispagestyle{plain}
{\sffamily\bfseries\fontsize{23}{26}\selectfont\color{ink}Где замыкается сольвус \(\alpha/\alpha_2\) в Ti--Al?\par}
{\large\sffamily\color{muted}Расчёт на MTP Qi при нулевом давлении\par}
\vspace{1mm}{\color{teal}\rule{\linewidth}{1.1pt}}
@FINDING@

\subsection*{Температура и область, охваченная прямым расчётом}
@SUMMARYTABLE@
@INTERPRETATION@
@RIGHTBRANCH@
@PROFILE@
\clearpage
\section*{Составы и барьер на конечных размерах}
@PLOTS@

При равных суммарных весах двух фаз химическая восприимчивость \(\partial\langle x\rangle/\partial\Delta\mu=\beta N\operatorname{Var}(x)\) содержит межфазный вклад \(\beta N(\Delta x)^2/4\), если \(\Delta x\) выражено в атомных долях. При сближении составов этот вклад исчезает. Продолжающийся в этой области разрыв порядка можно обнаружить по профилю \(Q\) и по энергии двух фаз.

При 1245~К в поле максимального отношения весов (\(q_c=0{,}47\)) барьер от меньшего пика растёт с 0,84 до 1,78 и 2,84~\(k_{\rm B}T\) для 768, 1400 и 2304 атомов. На 2304 атомах моды расположены при \(Q=0{,}083\) и 0,793 (рис.~\ref{fig:order}). При 1250~К на 768 атомах обе моды всё ещё наблюдаются, но максимальный вес одного варианта D0$_{19}$ составляет лишь 0,30 от веса разупорядоченного бассейна. Равновесная пара по отношению \(4:1\) в этой серии отсутствует.
@LIMITS@
\clearpage
\section*{Как определялись равновесные состояния}
Полубольшой изобарический Монте-Карло с \texttt{fitted.mtp} из набора Qi учитывает взаимные замены Ti и Al и классические ангармонические смещения. Относительные свободные энергии получены из равновесных вероятностей. В проходе выполняются по \(N\) проб замены и смещения, одна проба объёма, а раз в три прохода --- проба \(\log(c/a)\). В весе изменения объёма учитывается якобиан \((V'/V)^N\). Область ГПУ-кристалла задана условием \(W<0{,}40\). Здесь \(W=\max_{i,\mathbf d}2(\mathbf u_i-\bar{\mathbf u})\cdot\mathbf d/|\mathbf d|^2\), где \(\mathbf u_i\) --- смещение от узла, а \(\mathbf d\) пробегает двенадцать связей идеальной решётки. Допустимы объёмы 15--21~\AA$^3$/атом, \(1{,}45\le c/a\le1{,}72\); для проб формы и объёма требуется \(d_{\rm nn}^{\rm ideal}(1-W_{\max})\ge1{,}68052\)~\AA. Границы защищают от нефизического короткодействующего поведения MTP. \href{https://doi.org/10.1103/PhysRevMaterials.7.103602}{Qi и соавт. (2023)}, \href{https://doi.org/10.6084/m9.figshare.24079611}{потенциал и данные}.

В план расчётов входят ячейки из 96, 288, 768, 1400 и 2304 атомов. Соответствующие размеры в двухатомных ячейках ГПУ равны \(4{\times}4{\times}3\), \(6{\times}6{\times}4\), \(8{\times}8{\times}6\), \(10{\times}10{\times}7\) и \(12{\times}12{\times}8\). Через \(n_j\) обозначим число атомов Al на подрешётке \(j\), через \(k=\sum_jn_j\) их полное число. В каждое окно добавлен сдвиг энергии всей ячейки
\[
 U_{\rm bias}=\tfrac12 K(Q-q_0)^2-\Delta\mu\, k,\qquad
 Q=\frac{4(4\max_jn_j-k)}{3N}.
\]
Здесь \(\Delta\mu\) --- конфигурационное химическое поле на одну замену Ti на Al. Метод MBAR объединяет окна одного размера, температуры и механического ансамбля. Для проверки выборок учитываются временные корреляции, перекрытие соседних окон и вклад обеих фаз в оценку весов. \href{https://doi.org/10.1063/1.2978177}{Shirts и Chodera (2008).}

Поле \(\Delta\mu^*\) для составов выбирается по весам фаз. Четыре эквивалентных варианта D0$_{19}$ объединены операцией максимума в \(Q\), поэтому применяется условие \(P(\alpha_2)/P(\alpha)=4\). По разные стороны минимума между модами усредняется \(x=k/N\). Для оценки барьера поле отдельно подстраивается до равенства высот двух пиков плотности \(P(Q)\). Тогда \(\beta\Delta F=\ln(P_{\max}/P_{\min})\). Если седловая область между хорошо разделёнными фазами содержит две плоские границы при периодических условиях, главный вклад в барьер пропорционален их площади, порядка \(N^{2/3}\). Доменная геометрия, когерентные напряжения и корреляция \(Q\) со скрытыми степенями свободы могут изменить зависимость на малых ячейках. \href{https://doi.org/10.1103/PhysRevLett.65.137}{Lee и Kosterlitz (1990)}; \href{https://doi.org/10.1103/PhysRevB.47.14757}{Janke (1993).}

\subsection*{Конгруэнтное разупорядочение}
При фиксированном разделителе \(q_c\) обозначим \(R=P(Q>q_c)/[4P(Q\leq q_c)]\). Вершину области упорядочения ищем при температуре, где внутренний максимум \(\ln R\) по химическому полю равен нулю. На Ti-богатой ветви \(\ln R\) пересекает ноль снизу вверх с ростом поля, на Al-богатой --- сверху вниз. В точке слияния корней
\[
 \frac{\partial\ln R}{\partial\Delta\mu}=\beta N\,
      (x_{\alpha_2}-x_\alpha)=0.
\]
Совпадение составов в максимуме следует из этого уравнения. Независимые признаки разрыва --- разделённые моды порядка и разность энергий двух бассейнов. При одинаковом \(x\) и \(p=0\) оценка теплоты разупорядочения на атом принимает вид
\[
  \ell_{\rm dis} =
   \frac{\langle E-\Delta\mu\, k\rangle_\alpha -
         \langle E-\Delta\mu\, k\rangle_{\alpha_2}}{N}.
\]
Слагаемое \(\Delta\mu\,\Delta x\) исчезает, классические кинетические вклады фаз одинаковы. Положительное \(\ell_{\rm dis}\) соответствует поглощению теплоты. Модель концентрационных волн для ГПУ Ti--Al также допускает скачок порядка при стехиометрии. \href{https://arxiv.org/pdf/1406.0162}{Radchenko и соавт., рис.~4a.}

\subsection*{Переход к объёмному пределу}
Подгонка \(T_*(N)=T_\infty+a/N\) по 768, 1400 и 2304 атомам даёт 1242,5~К [1240,4; 1244,7]. Замена \(1/N\) на \(N^{-2/3}\) даёт 1242,4~К [1239,4; 1245,5]. По двум наибольшим ячейкам зависимость от \(1/N\) даёт 1246,5~К [1242,5; 1251,0]. Указанные 95\%-е блочные интервалы условны по форме каждой подгонки. Рабочая полоса 1240--1247~К характеризует межсерийный разброс и центральные экстраполяции; её статистическое покрытие неизвестно.
\clearpage
\section*{Независимые ориентиры и пределы вывода}
@THERMAL@
\subsection*{Что измеряет гистерезис}
Nichols и соавт. наблюдали самопроизвольный переход на MTP Qi по траекториям длиной 40~пс на 1024 атомах. Составы перед переходами из \(\alpha\) и D0$_{19}$ равны 15,2/23,4~ат.\% Al при 1000~К и 18,2/23,7 при 1100~К. При 1200~К в табл.~2 работы стоят прочерки. Наши равновесные составы при 1000~К равны 14,98/22,44~ат.\% Al. Границы гистерезиса зависят также от времени наблюдения и начального состояния. \href{https://doi.org/10.3390/met16020140}{Nichols и соавт. (2026), разд.~2.4, табл.~2.}

\subsection*{Кластерное разложение и эксперимент}
При 1200~К квантовый C46 на том же MTP даёт 23,46/24,91~ат.\% Al (\(N=5120\)), прямой MTP --- 22,00/24,31~ат.\% Al (\(N=1400\)). Ширина области составляет 1,44 и 2,31~п.п. соответственно. В квантовом C46 при 1225, 1250 и 1300~К отдельные контакты не разрешены. В классическом C46 при 1273,15~К химический потенциал монотонен от 13,98 до 36,86~ат.\% Al. В сравнение входят усечение кластерного базиса, квазигармоническое приближение и фиксация \(c/a\) в C46. Квантовая статистика колебаний также отличается от классической, использованной при прямом расчёте.

Для диффузионной пары, отожжённой при 1273,15~К в течение 200~ч, электронно-зондовый микроанализ дал 19,5\(\pm\)0,5 и 23,0\(\pm\)0,3~ат.\% Al. При 1323,15~К получены 21,3\(\pm\)0,3 и 23,5\(\pm\)0,3, а при 1423,15~К --- 24,7\(\pm\)0,2 и 25,7\(\pm\)0,2~ат.\% Al. Это измерения фазовых границ на указанных изотермах. Знаки \(\pm\) передают точность определения составов. \href{https://doi.org/10.1016/0966-9795(94)90018-3}{Kainuma, Palm и Inden (1994), табл.~1.}

Для CALPHAD-модели Liang составы фаз совпадают при 1437,28~К и 29,76~ат.\% Al. Эта расчётная вершина выше оценки прямого MTP примерно на 193~К и богаче алюминием на 2,6~п.п. Измеренная пара при 1423,15~К непосредственно показывает сохранение двухфазной области на 179~К выше вершины MTP. \href{https://doi.org/10.1007/s11669-024-01124-9}{Liang и соавт. (2024), табл.~3.}

{\footnotesize\color{muted}Проверены 36 заданий UMT, 972 траектории и 738\,700 продуктивных кадров. Численные данные сохранены в \texttt{analysis/closure\_summary.csv}, \texttt{analysis/critical\_order\_profiles.csv} и \texttt{analysis/thermal/}. Координаты графиков встроены в исходник PDF.}
\end{document}
